\section{Discussion}
\label{sec:discussion}
The calculations support a coherent picture: the backward channel is strongest near the edge and decays with distance, while the incident channel remains the dominant positive-power branch. The primary isolated backward-channel self-flux at \SI{2}{\micro\metre} is reported without claiming that it is total reflectance or that the coupling is experimentally calibrated. The matched-H comparison from the historical package separates edge conversion, \SI{2}{\micro\metre} return and \SI{100}{\micro\metre} return; those numbers are reproduced here only where supported by the same model.

The present package does not establish novelty over Qu et al.; the exact equivalence/novelty remains unresolved. Synthetic $A$, $f_{3333}$, and $H$ values are not presented as experimentally calibrated ZnO constants. The ordinary sourced constants are traceable to Zhang et al. and NIST SI definitions.

The main limitation is the model geometry: it is 2-D plane strain with a
specified propagation-plane tangential component, not a true 3-D oblique edge
solution. The open-circuit edge condition required an explicit gauge to remove
the electric-potential null space of the weak form; with that gauge the open
case is solvable and no longer fails on the conditioning of the tangent
matrix. The gauge removes a null space, not physics, and it is not claimed as
an independent validation of the open-boundary response.

The figures support three quantitative statements and nothing more. First, the
backward reflectance is governed almost entirely by the transverse wavenumber
and by the position of the operating point on the guided branch, with the
carrier--strain-gradient coupling acting as a perturbation below
\SI{2.5}{\percent}. Second, the conversion maximum coincides with the point
where the group velocity of the branch vanishes, so the effect is a
slow-mode effect rather than a resonance of the coupling term alone. Third, the
near field at the reference plane is reactive and must be excluded from any
transport argument, which is why the energy budget of Fig.~\ref{fig:fig05} is
closed on three channels and reports the forward-propagating channel as exactly
zero.
